\chapter{Supplementary vision experiments}
\label{app:supplementary-vision-experiments}

This appendix keeps the vision-track details that are useful for reproducibility
and auditability but are too detailed for the main argument. The main empirical
claims are presented in Chapter~\ref{chap:vision-curriculum}; this appendix
records the protocol choices and secondary diagnostics behind them.

\section{CIFAR-100 model-size study}
\label{app:cifar100-model-size-curriculum}

The screening study used CIFAR-100 with one baseline and five fixed curriculum
lengths for each model configuration. Most runs used seed 42, with repeated
baseline runs averaged where available. The common training budget was 100
epochs, Adam with learning rate 0.001, and curriculum lengths of 5, 10, 20, 30,
and 40 epochs.

\begin{table}[h]
\centering
\small
\begin{tabular}{lrrrrrr}
\toprule
Model & Base & Curr5 & Curr10 & Curr20 & Curr30 & Curr40 \\
\midrule
\texttt{cnn\_w0.5} & 34.78 & 37.57 & 37.83 & 37.96 & 37.29 & 36.34 \\
\texttt{cnn\_w1} & 38.49 & 41.88 & 42.11 & 40.70 & 38.20 & 37.08 \\
\texttt{cnn\_w2} & 42.56 & 44.72 & 42.50 & 40.30 & 38.24 & 38.28 \\
\texttt{cnn\_w4} & 44.63 & 44.47 & 43.19 & 40.41 & 39.50 & 38.25 \\
\texttt{resnet18} & 68.76 & 68.84 & 69.15 & 69.37 & 68.49 & 68.71 \\
\bottomrule
\end{tabular}
\caption{Selected CIFAR-100 curriculum-length results. Values are best test
accuracies in percent. The table supports the main-text claim that the best
curriculum length shortens as model capacity grows.}
\label{tab:app-cifar100-lengths-compact}
\end{table}

Secondary metrics such as macro-F1, top-5 accuracy, calibration error, and
hierarchy-aware scores generally followed the accuracy result only loosely.
Therefore the main thesis treats top-1 accuracy as the primary measure and uses
secondary metrics only as diagnostics.

\section{Adaptive schedules, dataset controls, and roughness}
\label{app:roughness-followup-curriculum}

The follow-up repeated important CIFAR-100 configurations over seeds 42, 43, and
44 and added adaptive plateau schedules. The adaptive rule advanced from a
coarse stage when fine-label validation accuracy stopped improving after a
minimum stage length or when a maximum stage length was reached. Two hierarchy
coverage variants were used. The all-level variant could visit every generated
non-singleton hierarchy level. The first-level variant used only the first
non-singleton level, meaning the broadest usable coarse partition after removing
the trivial root cluster, and then switched to fine labels. This rule was stable
to run, but it often spent too many epochs in coarse stages.

Dataset controls used CIFAR-10, Fashion-MNIST, and Tiny ImageNet. CIFAR-10
showed smaller gains than CIFAR-100, Fashion-MNIST was mostly neutral or
negative, and Tiny ImageNet depended strongly on model and optimizer setup. The
roughness and sharpness probes were treated as diagnostics only. They sometimes
supported a flatter-minima interpretation, but the pattern was not consistent
enough to establish a mechanism.

\begin{table}[h]
\centering
\small
\begin{tabular}{lll}
\toprule
Diagnostic & Main observation & Thesis use \\
\midrule
Adaptive schedule & often too much coarse training & explains best-fixed advantage here \\
Dataset controls & CIFAR-100 strongest & shows structure matters \\
Trajectory AUC & usually negative & peak gain is not faster training \\
Sharpness/Hessian & mixed signs & no strong mechanism claim \\
\bottomrule
\end{tabular}
\caption{Supplementary interpretation of the vision follow-up diagnostics.}
\label{tab:app-vision-diagnostic-summary}
\end{table}

\section{Random-hierarchy ablation}
\label{app:hierarchy-ablation-curriculum}

The random-hierarchy ablation preserved the learned hierarchy shape and cluster
sizes but randomly reassigned CIFAR-100 classes to clusters. For each training
seed, three random hierarchy seeds were evaluated. This made the random control
stronger than simply comparing against an arbitrary hand-made hierarchy: the
amount of label coarsening was held fixed, and only the grouping information was
destroyed.

The learned hierarchy outperformed the random mean for all three training seeds.
At the same time, the random mean was above the direct baseline, which is why the main
text separates a general coarse-label warm-up effect from the additional benefit
of meaningful hierarchy structure.
